\chapter{Literature Review}\label{ch:literature}

\section{Related Work}\label{sec:related}

\shap{} assigns each feature the Shapley value of its marginal contribution, with
\textsc{TreeSHAP} making this tractable for tree
ensembles~\cite{lundberg2017unified,lundberg2020local}. \lime{} perturbs the instance,
weights the perturbations by proximity, and fits a sparse linear surrogate whose
coefficients are the explanation~\cite{ribeiro2016why}. The difference is not cosmetic:
one attributes the model's output, the other approximates its local function.

Because there is no ground-truth attribution, evaluation is indirect. Perturbation tests
delete features in explanation order and measure how fast the prediction degrades,
summarised as area over the perturbation curve~\cite{samek2017evaluating}; \textsc{eraser}
formalised comprehensiveness and sufficiency~\cite{deyoung2020eraser}. Hooker et
al.~\cite{hooker2019benchmark} showed such tests are confounded by the off-manifold
inputs perturbation creates, which we address directly in Section~\ref{sec:pitfalls}.
Robustness is usually estimated as a local Lipschitz
constant~\cite{alvarezmelis2018robustness}, and rank-biased overlap compares ranked lists
without requiring a cut-off~\cite{webber2010similarity}. Neither method is as strong in
practice as in theory --- the background distribution changes \shap{}'s attributions,
which Section~\ref{sec:pitfalls} treats as a design decision --- and both can be
deliberately fooled~\cite{slack2020fooling}. We take these as reasons for careful measurement rather
than settled verdicts.

Combining several explainers has been proposed before, both as an ensemble of
attributions and as a framework for aggregating explanation
outputs~\cite{bhatt2020aggregating}; rank fusion itself is standard in information
retrieval, where reciprocal-rank fusion is a strong and very simple
baseline~\cite{cormack2009reciprocal}. We therefore claim no priority for the idea of
aggregation, and use both of those rules as baselines. What is new here is weighting the
two explainers by their \emph{measured} fidelity rather than equally, fitting that weight
on a disjoint calibration split so the rule cannot be tuned on its own test set, and
testing the result against the better of the two rather than their mean.

Several studies contrast \shap{} and \lime{} on tabular and clinical
data~\cite{salih2024perspective,duell2021comparison}. Most vary a single axis, none we
surveyed reports a random-ranking control or states its attribution space, and none
pre-registers its hypotheses.


\section{Synthesis: what this literature leaves open}
Three gaps recur across the work surveyed above, and each maps to a design decision in
Chapter~\ref{ch:method}.

\textbf{Scope.} Most comparisons vary one axis --- several models on one dataset, or
several datasets with one model. Because explanation behaviour plausibly depends on
both the model's functional form and the data's feature redundancy, a single-axis
result cannot separate the two. We vary both, giving eighteen cells.

\textbf{Inference.} Point differences are reported without hypothesis tests, effect
sizes, or correction for multiplicity. With per-instance metrics running to hundreds of
observations per cell, almost any difference is `significant', so the absence of effect
sizes is the more serious omission. We report effect sizes and intervals on every
comparison and correct within hypothesis family.

\textbf{Comparability.} No surveyed comparison states which output space its
attributions live in. Since LIME explains predicted probability while tree SHAP
defaults to a log-odds margin, a comparison that does not fix this is contrasting two
different quantities. We fix it and verify it (Section~\ref{sec:pitfalls}).

A fourth observation concerns the evaluation metrics themselves. Fidelity is reported
without a random baseline in every comparison we surveyed, which leaves the absolute
numbers uninterpretable, and stability is reported as a single quantity even though at
least three distinct notions are in circulation. Both shape our design.
